\documentclass[aspectratio=169]{beamer}
\usetheme{simple}

\input{preamble/preamble.tex}
\input{preamble/preamble_math.tex}
% \input{preamble/preamble_acronyms.tex}

\title{Variational Autoencoders}
\subtitle{STATS305B: Applied  Statistics II}
\author{Scott Linderman}
\date{\today}


\begin{document}


\maketitle

\begin{frame}{Last Time...}
    \textbf{Outline:}
    \begin{itemize}
        \item Principal Components Analysis (PCA)
        \item PCA as a linear Gaussian latent variable model
        \item Factor analysis
        \item Linear Dynamical Systems \& the Kalman Filter/Smoother
    \end{itemize}
    \vspace{1em}
\end{frame}

\begin{frame}{Today...}
    \textbf{Outline:}
    \begin{itemize}
        \item (Probabilistic) PCA as a (stochastic) linear autoencoder
        \item Variational Autoencoders (VAEs)
        \item Review of Gradient-based Variational Inference
        \item Demo
    \end{itemize}
    \vspace{1em}
\end{frame}

   
\begin{frame}[t,allowframebreaks]{PCA as a Linear
Autoencoder}\label{pca-as-a-linear-autoencoder}

Last time we introduced PCA as a method for finding dimensions of
maximal variance. However, we can arrive at the same model from another
perspective: find the linear projection that minimizes the average
projection cost.

To formalize this, let \begin{align*}
    \mbW = 
    \begin{bmatrix}
    | & & | \\
    \mbw_1 & \cdots & \mbw_M \\
    | & & |
    \end{bmatrix} \in \reals^{D \times M} \quad \text{with} \quad \mbW^\top \mbW = \mbI
\end{align*} be an orthogonal basis for the principal subspace.

We will \textbf{encode} each data point by subtracting the mean and
projecting onto the principal subspace to obtain
\(\mbz_n = \mbW^\top (\mbx_n - \bar{\mbx})\).

Since \(\mbW\) is an orthogonal matrix, all we need to do to
\textbf{decode} the encoded data point is multiply \(\mbW \mbz_n\) and
add back the mean. That gives us, \begin{align*}
    \hat{\mbx}_n = \mbW \mbz_n + \bar{\mbx} = \mbW \mbW^\top (\mbx_n - \bar{\mbx}) + \bar{\mbx}.
\end{align*}

\textbf{Goal:} Find an orthogonal matrix \(\mbW\) that \emph{minimizes}
the mean squared reconstruction error, \begin{align*}
    \cL(\mbW) &= \frac{1}{N} \sum_{n=1}^N \|\mbx_n - \hat{\mbx}_n\|_2^2
\end{align*} We can write this in matrix notation instead. Let
\(\mbX \in \reals^{N \times D}\) be the \emph{centered data matrix} with
rows \((\mbx_n - \bar{\mbx})^\top\). Then, \begin{align*}
    \cL(\mbW) 
    &= \frac{1}{N} \Tr[(\mbX - \hat{\mbX})^\top (\mbX - \hat{\mbX})] \\
    &= \frac{1}{N} \Tr[(\mbX - \mbX \mbW \mbW^\top)^\top (\mbX - \mbX \mbW \mbW^\top)] \\
    &= \frac{1}{N} \Tr[(\mbX(\mbI - \mbW \mbW^\top))^\top (\mbX (\mbI - \mbW \mbW^\top))]\\
    &= \frac{1}{N} \Tr[(\mbI - \mbW \mbW^\top)^\top \mbX^\top \mbX (\mbI - \mbW \mbW^\top)] \\
    &= \Tr[(\mbI - \mbW \mbW^\top)^\top \mbS (\mbI - \mbW \mbW^\top)]
\end{align*} where \(\mbS = \frac{1}{N} \mbX^\top \mbX\) is the sample
covariance matrix.

Now apply the circular trace property, 
\begin{align*}
    \cL(\mbW) 
    &= \Tr[\mbS (\mbI - \mbW \mbW^\top) (\mbI - \mbW \mbW^\top)^\top]
\end{align*}

\textbf{Question:} What does
\((\mbI - \mbW \mbW^\top) (\mbI - \mbW \mbW^\top)^\top\) equal? 

Note
that \(\mbI - \mbW \mbW^\top\) is a projection matrix --- it projects a
vector onto the \emph{nullspace} of \(\mbW\). Applying the projection
operator twice doesn't change the result. Mathematically, 
\begin{align*}
    (\mbI - \mbW \mbW^\top) (\mbI - \mbW \mbW^\top)^\top 
    = \mbI - 2 \mbW \mbW^\top + \mbW \mbW^\top \mbW \mbW^\top
    = \mbI - \mbW \mbW^\top
\end{align*} 
where we used the fact that \(\mbW^\top \mbW = \mbI\) since
\(\mbW\) is an orthogonal matrix. 

Thus, the objective simplifies to, \begin{align*}
    \cL(\mbW) 
    = \Tr[\mbS(\mbI - \mbW \mbW^\top)]
    = \Tr[\mbS] - \Tr[\mbS \mbW \mbW^\top]
    = \mathrm{const} - \Tr[\mbW^\top \mbS \mbW].
\end{align*} Let \(\mbU \mbLambda \mbU^\top\) be the eigendecomposition
of \(\mbS\). (Since it is a covariance matrix, the eigenvectors are
orthogonal.) Plugging in, \begin{align*}
    \cL(\mbW) 
    &= \mathrm{const} - \Tr[\mbW^\top \mbU \mbLambda \mbU^\top \mbW] \\
    &= \mathrm{const} - \Tr\left[\mbW^\top \left( \sum_{d=1}^D \lambda_d \mbu_d \mbu_d^\top \right) \mbW\right] \\
    &= \mathrm{const} - \sum_{m=1}^M \sum_{d=1}^D \lambda_d \mbw_m^\top \mbu_d \mbu_d^\top \mbw_m \\
    &= \mathrm{const} - \sum_{m=1}^M \sum_{d=1}^D \lambda_d (\mbw_m^\top \mbu_d)^2
\end{align*} We want to minimize \(\cL(\mbW)\) subject to \(\mbW\) being
orthogonal. What is the solution? \(\mbW = \mbU_M\)!
\end{frame}

\begin{frame}[t,allowframebreaks]{Probabilistic PCA as a Stochastic Linear
Autoencoder}\label{probabilistic-pca-as-a-stochastic-linear-autoencoder}

The jump from PCA to probabilistic PCA was to treat the scores,
\(\mbz_n\), as latent variables. We gave them a Gaussian prior and
assumed a generative model, 
\begin{align*}
    \mbz_n &\iid{\sim} \mathrm{N}(\mbzero, \mbI) \\
    \mbx_n \mid \mbz_n &\sim \mathrm{N}(\mbW \mbz + \mbmu, \sigma^2 \mbI),
\end{align*} 
where \(\mbz_n \in \reals^M\) is a latent variable,
\(\mbW \in \reals^{D \times M}\) are the weights, \(\mbmu \in \reals^D\)
is the bias parameter, and \(\sigma^2 \in \reals_+\) is a variance.

Given the parameters, we showed that the posterior distribution over
latent variables is, 
\begin{align*}
p(\mbz_n \mid \mbx_n; \mbtheta)
&=
\mathrm{N}(\mbz_n \mid (\sigma^2 \mbI + \mbW^\top \mbW)^{-1} \mbW^\top (\mbx_n - \mbmu), \, \sigma^2 (\sigma^2 \mbI + \mbW^\top \mbW)^{-1}),
\end{align*} 
and when \(\sigma^2 \to 0\) and \(\mbW\) are the PCs, this
converges to a delta function on the PCA scores.

To simplify notation, let's rewrite the posterior as, 
\begin{align*}
p(\mbz_n \mid \mbx_n; \mbtheta)
&= \mathrm{N}(\mbz_n \mid \mbE \mbx_n + \mbd, \mbG)
\end{align*} 
where 
\begin{align*}
\mbE &= (\sigma^2 \mbI + \mbW^\top \mbW)^{-1} \mbW^\top \\
\mbd &= - (\sigma^2 \mbI + \mbW^\top \mbW)^{-1} \mbW^\top \mbmu \\
\mbG &= \sigma^2 (\sigma^2 \mbI + \mbW^\top \mbW)^{-1}.
\end{align*}

To ``autoencode'' the data using probabilistic PCA, we could sample
\(\mbz_n \sim p(\mbz_n \mid \mbx_n; \mbtheta)\) by mapping
\(\mbx_n \mapsto \mbE \mbx_n + \mbd\) and adding Gaussian noise with
covariance \(\mbG\). Then we can generate a new data point
\(\hat{\mbx}_n \sim \mathrm{N}(\mbW \mbz_n + \mbmu, \sigma^2 \mbI)\).

Just as PCA could be motivated as finding weights that minimize the
reconstruction error, probabilistic PCA can be seen as finding weights
that minimize the \emph{expected} reconstruction error (subject to a
little regularization).
\end{frame}

\begin{frame}[t,allowframebreaks]{Revisiting EM for Probabilistic
PCA}\label{revisiting-em-for-probabilistic-pca}

Recall that we justified EM by connecting it to variational inference.
We showed that EM maximizes the evidence lower bound (ELBO),
\begin{align*}
\cL(q, \mbtheta) 
&= \sum_n \mathbb{E}_{q(\mbz_n)}\left[\log p(\mbx_n, \mbz_n; \mbtheta) - \log q(\mbz_n) \right] \\
&= \sum_n \underbrace{\mathbb{E}_{q(\mbz_n)}\left[\log p(\mbx_n \mid \mbz_n; \mbtheta)\right]}_{\text{expected log likelihood}} - \underbrace{\KL{q(\mbz_n)}{p(\mbz_n)}}_{\text{KL to the prior}} \\
&\leq \sum_n \log p(\mbx_n; \mbtheta)
\end{align*} 
EM can be seen as coordinate ascent on the ELBO: 
\begin{enumerate}
    \item The E-step sets \(q\) to the posterior over latent variables, \(q(\mbz_n) = p(\mbz_n \mid \mbx_n; \mbtheta)\) for the current parameters \(\mbtheta\). 
    \item 2. The M-step updates the parameters to maximize the expected log joint probability.
\end{enumerate}

In probabilistic PCA, maximizing the expected log likelihood is the same
as minimizing the expected reconstruction error. Plugging in the
definition of the model and the posterior, 
\begin{align*}
\mathbb{E}_{q(\mbz_n)}\left[\log p(\mbx_n \mid \mbz_n; \mbtheta)\right]
&= \E_{\mathrm{N}(\mbz_n \mid \mbE \mbx_n + \mbd, \mbG)}\left[\log \mathrm{N}(\mbx_n \mid \mbW \mbz_n + \mbmu, \sigma^2 \mbI) \right] \\
&= -\frac{1}{2\sigma^2 }\E_{\mathrm{N}(\mbz_n \mid \mbE \mbx_n + \mbd, \mbG)}\left[\|\mbx_n - \mbW \mbz_n - \mbmu \|_2^2 \right] + \log \sigma + c \\
\end{align*} 
where \(c\) is an additive constant, and \(\mbE\),
\(\mbd\), \(\mbG\) are functions of the parameters, as defined above,

We can't forget about the KL to the prior though! When we maximize the
ELBO, we are essentially finding weights that minimize the expected
reconstruction error, while also not deviating too far from the standard
normal prior on the latent variables.
\end{frame}

\begin{frame}[t,allowframebreaks]{Variational Autoencoders}\label{variational-autoencoders}

Variational Autoencodres (VAEs) are ``deep'' but conceptually simple
generative models. The generative model is the same as probabilistic PCA, but allowing for \textbf{nonlinear} mappings between latent variables and
observations.

To sample a data point \(\mbx_n\), 
\begin{enumerate}
    \item First, sample \textbf{latent variables} \(\mbz_n\), 
    \begin{align*}
    \mbz_n &\sim \mathrm{N}(\mbzero, \mbI) 
    \end{align*}

    \item Then sample the data point \(\mbx_n\) from a conditional distribution with mean, 
    \begin{align*}
    \E[\mbx_n \mid \mbz_n] &= g(\mbz_n; \mbtheta),
    \end{align*} 
    where \(g: \reals^H \to \reals^D\) is a nonlinear mapping parameterized by \(\mbtheta\).
\end{enumerate}

We will assume \(g\) is a simple \textbf{feedforward neural network} of
the form, 
\begin{align*}
    g(\mbz; \mbtheta) &= g_L(g_{L-1}(\cdots g_1(\mbz) \cdots))
\end{align*} 
where each \textbf{layer} is a cascade of a linear mapping
followed by an element-wise nonlinearity (except for the last layer,
perhaps). For example, 
\begin{align*}
    g_\ell(\mbu_{\ell}) = \mathrm{relu}(\mbW_\ell \mbu_{\ell} + \mbb_\ell); \qquad \mathrm{relu}(a) = \max(0, a).
\end{align*} 
The generative parameters consist of the weights and
biases, \(\mbtheta = \{\mbW_\ell, \mbb_\ell\}_{\ell=1}^L\).
\end{frame}

\begin{frame}[t,allowframebreaks]{Learning and Inference}\label{learning-and-inference}

We have two goals. The \textbf{learning goal} is to find the parameters
that \textbf{maximize the marginal likelihood of the data},
\begin{align*}
    \mbtheta^\star 
    &= \arg \max_{\mbtheta} p(\mbX; \mbtheta)\\
    &= \arg \max_{\mbtheta} \prod_{n=1}^N \int p(\mbx_n \mid \mbz_n; \mbtheta) \, p(\mbz_n; \mbtheta) \dif \mbz_n
\end{align*}

The \textbf{inference goal} is to find the \textbf{posterior
distribution of latent variables}, \begin{align*}
    p(\mbz_n \mid \mbx_n; \mbtheta) 
    &= \frac{p(\mbx_n \mid \mbz_n; \mbtheta) \, p(\mbz_n; \mbtheta)}{\int p(\mbx_n \mid \mbz_n'; \mbtheta) \, p(\mbz_n'; \mbtheta)\dif \mbz_n'}
\end{align*}

Both goals require an integral over \(\mbz_n\), but that is intractable
for deep generative models.
\end{frame}

\begin{frame}[t,allowframebreaks]{The Evidence Lower Bound
(ELBO)}\label{the-evidence-lower-bound-elbo}

\textbf{Idea:} Use the ELBO to get a bound on the marginal probability
and maximize that instead. \begin{align*}
\log p(\mbX ; \mbtheta) 
&= \sum_{n=1}^N \log p(\mbx_n; \mbtheta) \\
&\geq \sum_{n=1}^N \log p(\mbx_n; \mbtheta) - \KL{q(\mbz_n; \mblambda_n)}{p(\mbz_n \mid \mbx_n; \mbtheta)} \\
&= \sum_{n=1}^N \underbrace{\E_{q(\mbz_n)}\left[ \log p(\mbx_n, \mbz_n; \mbtheta) - \log q(\mbz_n; \mblambda_n) \right]}_{\text{"local ELBO"}} \\
&\triangleq \sum_{n=1}^N \cL_n(\mblambda_n, \mbtheta) 
= \cL(\mblambda, \mbtheta)
\end{align*} where \(\mblambda = \{\mblambda_n\}_{n=1}^N\).
Here, I've written the ELBO as a sum of \emph{local ELBOs} \(\cL_n\).
\end{frame}

\begin{frame}[t,allowframebreaks]{Variational Inference}\label{variational-inference}

The ELBO is still maximized (and the bound is tight) when each \(q\) is
equal to the true posterior, \begin{align*}
    q(\mbz_n; \mblambda_n) &= p(\mbz_n \mid \mbx_n, \mbtheta).
\end{align*} Unfortunately, the posterior no longer has a simple, closed
form.

\textbf{Question:} Suppose
\(\mbx_n \sim \mathrm{N}(g(\mbz_n; \mbtheta), \mbI)\). This deep
generative model has a Gaussian prior on \(\mbz_n\) and a Gaussian
likelihood for \(\mbx_n\) given \(\mbz_n\). Why isn't the posterior
Gaussian? 

Nevertheless, we can still constrain \(q\) to belong to a simple family.
For example, we could constrain it to be Gaussian and seek the best
Gaussian approximation to the posterior. This is sometimes called
\textbf{fixed-form variational inference}. Let, \begin{align*}
    \cQ = \left\{q: q(\mbz; \mblambda) = \mathrm{N}\big(\mbz \mid \mbmu, \diag(\mbsigma^2)\big) \text{ for } \mblambda = (\mbmu, \log \mbsigma^2) \in \reals^{2H} \right\}
\end{align*}

Then, for fixed parameters \(\mbtheta\), the best \(q\) in this
\textbf{variational family} is, \begin{align*}
    q^\star 
    &= \arg \min_{q \in \cQ} \KL{q(\mbz_n; \mblambda_n)}{p(\mbz_n \mid \mbx_n; \mbtheta)} \\
    &= \arg \max_{\mblambda_n \in \reals^{2H}} \cL_n(\mblambda_n, \mbtheta).
\end{align*}
\end{frame}

\begin{frame}[t,allowframebreaks]{Variational Expectation-Maximization
(vEM)}\label{variational-expectation-maximization-vem}

Now we can introduce a new algorithm.

\textbf{Algorithm"} Variational EM (vEM) Repeat until either the ELBO
or the parameters converges: 
\begin{enumerate}
    \item \textbf{M-step:} Set 
    \(\mbtheta \leftarrow \arg \max_{\mbtheta} \cL(\mblambda, \mbtheta)\) 
    \item \textbf{E-step:} Set 
    \(\mblambda_n \leftarrow \arg \max_{\mblambda_n \in \mbLambda} \cL_n(\mblambda_n, \mbtheta)\)
    for \(n=1,\ldots,N\) 
    \item Compute (an estimate of) the ELBO \(\cL(\mblambda, \mbtheta)\).
\end{enumerate}

In general, none of these steps will have closed form solutions, so
we'll have to use approximations.
\end{frame}

\begin{frame}[t,allowframebreaks]{Generic M-step with Stochastic Gradient
Ascent}\label{generic-m-step-with-stochastic-gradient-ascent}

For exponential family mixture models, the M-step had a closed form
solution. For deep generative models, we need a more general approach.

If the parameters are unconstrained and the ELBO is differentiable wrt
\(\mbtheta\), we can use stochastic gradient ascent. \begin{align*}
    \mbtheta &\leftarrow \mbtheta + \alpha \nabla_{\mbtheta} \cL(q, \mbtheta) \\
    &= \mbtheta + \alpha \sum_{n=1}^N \mathbb{E}_{q(\mbz_n; \mblambda_n)} \left[ \nabla_{\mbtheta} \log p(\mbx_n, \mbz_n; \mbtheta) \right]
\end{align*}

Note that the expected gradient wrt \(\mbtheta\) can be computed using
ordinary Monte Carlo --- nothing fancy needed!
\end{frame}

\begin{frame}[t,allowframebreaks]{The Variational E-step}\label{the-variational-e-step}

Assume \(\cQ\) is the family of Gaussian distributions with diagonal
covariance: 
\begin{align*}
    \cQ = \left\{q: q(\mbz; \mblambda) = \mathrm{N}\big(\mbz \mid \mbmu, \diag(\mbsigma^2)\big) \text{ for } \mblambda = (\mbmu, \log \mbsigma^2) \in \reals^{2H} \right\}
\end{align*} 
This family is indexed by \textbf{variational parameters}
\(\mblambda_n = (\mbmu_n, \log \mbsigma_n^2) \in \reals^{2H}\).

To perform SGD, we need an unbiased estimate of the gradient of the
local ELBO, but \begin{align*}
\nabla_{\mblambda_n} \cL_n(\mblambda_n, \mbtheta) 
&= \nabla_{\mblambda_n} \E_{q(\mbz_n; \mblambda_n)} \left[ \log p(\mbx_n, \mbz_n; \mbtheta) - \log q(\mbz_n; \mblambda_n) \right] \\
&\textcolor{red}{\neq} \;  \E_{q(\mbz_n; \mblambda_n)} \left[ \nabla_{\mblambda_n} \left(\log p(\mbx_n, \mbz_n; \mbtheta) - \log q(\mbz_n; \mblambda_n)\right) \right].
\end{align*}
\end{frame}

\begin{frame}[t,allowframebreaks]{Reparameterization Trick}\label{reparameterization-trick}

One way around this problem is to use the \textbf{reparameterization
trick}, aka the \textbf{pathwise gradient estimator}. Note that,
\begin{align*}
    \mbz_n \sim q(\mbz_n; \mblambda_n) \quad \iff \quad
    \mbz_n &= r(\mblambda_n, \mbepsilon), \quad \mbepsilon \sim \mathrm{N}(\mbzero, \mbI) 
\end{align*} where
\(r(\mblambda_n, \mbepsilon) = \mbmu_n + \mbsigma_n \mbepsilon\) is a
reparameterization of \(\mbz_n\) in terms of parameters \(\mblambda_n\)
and noise \(\mbepsilon\).

We can use the \textbf{law of the unconscious statistician} to rewrite
the expectations as, 
\begin{align*}
    \E_{q(\mbz_n; \mblambda_n)} \left[h(\mbx_n, \mbz_n, \mbtheta, \mblambda_n) \right]
    &= \E_{\mbepsilon \sim \mathrm{N}(\mbzero, \mbI)} \left[h(\mbx_n, r(\mblambda_n, \mbepsilon), \mbtheta, \mblambda_n) \right]
\end{align*} 
where 
\begin{align*}
h(\mbx_n, \mbz_n, \mbtheta, \mblambda_n) = \log p(\mbx_n, \mbz_n; \mbtheta) - \log q(\mbz_n; \mblambda_n).
\end{align*} 
The distribution that the expectation is taken under no
longer depends on the parameters \(\mblambda_n\), so we can simply take
the gradient inside the expectation, \begin{align*}
    \nabla_{\mblambda} \E_{q(\mbz_n; \mblambda_n)} \left[h(\mbx_n, \mbz_n, \mbtheta, \mblambda_n) \right]
    &=  \E_{\mbepsilon \sim \mathrm{N}(\mbzero, \mbI)} \left[\nabla_{\mblambda_n} h(\mbx_n, r(\mblambda_n, \mbepsilon), \mbtheta, \mblambda_n) \right]
\end{align*} 
Now we can use Monte Carlo to obtain an unbiased estimate
of the final expectation!
\end{frame}

\begin{frame}[t,allowframebreaks]{Working with mini-batches of
data}\label{working-with-mini-batches-of-data}

We can view the ELBO as an expectation over data indices, 
\begin{align*}
    \cL(\mblambda, \mbtheta) 
    &= \sum_{n=1}^N \cL_n(\mblambda_n, \mbtheta) \\
    &= N \, \E_{n \sim \mathrm{Unif}([N])}[\cL_n(\mblambda_n, \mbtheta)].
\end{align*} 
We can use Monte Carlo to approximate the expectation (and
its gradient) by drawing \textbf{mini-batches} of data points at random.

In practice, we often cycle through mini-batches of data points
deterministically. Each pass over the whole dataset is called an
\textbf{epoch.}
\end{frame}

\begin{frame}[t,allowframebreaks]{Algorithm}\label{algorithm}

Now we can add some detail to our variational expectation maximization
algorithm.

For epoch \(i=1,\ldots,\infty\):

For \(n=1,\ldots,N\):

\begin{enumerate}
\def\labelenumi{\arabic{enumi}.}
\item
  Sample \(\epsilon_n^{(m)} \iid{\sim} \mathrm{N}(\mbzero, \mbI)\) for
  \(m=1,\ldots,M\).
\item
  \textbf{M-Step}:

  \begin{enumerate}
  \def\labelenumii{\alph{enumii}.}
  \item
    Estimate \begin{align*}
     \hat{\nabla}_{\mbtheta} \cL_n(\mblambda_n, \mbtheta)
     &= 
     \frac{1}{M} \sum_{m=1}^M \left[ \nabla_{\mbtheta} \log p(\mbx_n, r(\mblambda_n, \mbepsilon_n^{(m)}); \mbtheta) \right]
     \end{align*}
  \item
    Set
    \(\mbtheta \leftarrow \mbtheta + \alpha_i N \hat{\nabla}_{\mbtheta} \cL_n(\mblambda_n, \mbtheta)\)
  \end{enumerate}
\item
  \textbf{E-step:}

  \begin{enumerate}
  \def\labelenumii{\alph{enumii}.}
  \item
    Estimate \begin{align*}
     \hat{\nabla}_{\mblambda} \cL_n(\mblambda_n, \mbtheta) 
     &= \frac{1}{M} \sum_{m=1}^M \nabla_{\mblambda} \left[\log p(\mbx_n, r(\mblambda_n, \mbepsilon_n^{(m)}); \mbtheta) - \log q(r(\mblambda_n, \mbepsilon_n^{(m)}), \mblambda_n) \right]
     \end{align*}
  \item
    Set
    \(\mblambda_n \leftarrow \mblambda_n + \alpha_i \hat{\nabla}_{\mblambda} \cL_n(\mblambda_n, \mbtheta)\).
  \end{enumerate}
\item
  Estimate the ELBO \begin{align*}
   \hat{\cL}(\mblambda, \mbtheta) 
   &= \frac{N}{M} \sum_{m=1}^M \log p(\mbx_n, r(\mblambda_n, \mbepsilon_n^{(m)}); \mbtheta) - \log q(r(\mblambda_n, \mbepsilon_n^{(m)}); \mblambda_n)
   \end{align*}
\item
  Decay step size \(\alpha_i\) according to schedule. 
\end{enumerate}
\end{frame}

\begin{frame}[t,allowframebreaks]{Amortized Inference}\label{amortized-inference}

Note that vEM involves optimizing separate variational parameters
\(\mblambda_n\) for each data point. For large datasets where we are
optimizing using mini-batches of data points, this leads to a strange
asymmetry: we update the generative model parameters \(\mbtheta\) every
mini-batch, but we only update the variational parameters for the
\(n\)-th data point once per epoch. Is there any way to share
information across data points?

Note that the optimal variational parameters are just a function of the
data point and the model parameters, 
\begin{align*}
    \mblambda_n^\star &= \arg \min_{\mblambda_n} \KL{q(\mbz_n; \mblambda_n)}{p(\mbz_n \mid \mbx_n, \mbtheta)} 
    \triangleq f^\star(\mbx_n, \mbtheta).
\end{align*} 
for some implicit and generally nonlinear function
\(f^\star\).

VAEs learn an approximation to \(f^\star(\mbx_n, \mbtheta)\) with an
\textbf{inference network}, a.k.a. \textbf{recognition network} or
\textbf{encoder}.

The inference network is (yet another) neural network that takes in a
data point \(\mbx_n\) and outputs variational parameters
\(\mblambda_n\), \begin{align*}
    \mblambda_n & \approx f(\mbx_n, \mbphi),
\end{align*} where \(\mbphi\) are the weights of the network.

The advantage is that the inference network shares information across
data points --- it \emph{amortizes} the cost of inference, hence the
name. The disadvantage is the output will not minimize the KL
divergence. However, in practice we might tolerate a worse variational
posterior and a weaker lower bound if it leads to faster optimization of
the ELBO overall.
\end{frame}

\begin{frame}[t,allowframebreaks]{Putting it all together}\label{putting-it-all-together}

Logically, I find it helpful to distinguish between the E and M steps,
but with recognition networks and stochastic gradient ascent, the line
is blurred.

The final algorithm looks like this.

\textbf{Variational EM (with amortized inference)}

Repeat until either the ELBO or the parameters converges:

\begin{enumerate}
\def\labelenumi{\arabic{enumi}.}
\item
  Sample data point \(n \sim \mathrm{Unif}(1, \ldots, N)\). {[}Or a
  minibatch of data points.{]}
\item
  Estimate the local ELBO \(\cL_n(\mbphi, \mbtheta)\) with Monte Carlo.
  {[}Note: it is a function of \(\mbphi\) instead of \(\mblambda_n\).{]}
\item
  Compute unbiased Monte Carlo estimates of the gradients
  \(\widehat{\nabla}_{\mbtheta} \cL_n(\mbphi, \mbtheta)\) and
  \(\widehat{\nabla}_{\mbphi} \cL_n(\mbphi, \mbtheta)\). {[}The latter
  requires the reparameterization trick.{]}
\item
  Set \begin{align*}
   \mbtheta &\leftarrow \mbtheta + \alpha_i \widehat{\nabla}_{\mbtheta} \cL_n(\mbphi, \mbtheta) \\
   \mbphi &\leftarrow \mbphi + \alpha_i \widehat{\nabla}_{\mbphi} \cL_n(\mbphi, \mbtheta)
  \end{align*} with step size \(\alpha_i\) decreasing over iterations
  \(i\) according to a valid schedule.
\end{enumerate}
\end{frame}

   
\end{document}
